\IfFileExists{stacks-project.cls}{%
\documentclass{stacks-project}
}{%
\documentclass{amsart}
}

% For dealing with references we use the comment environment
\usepackage{verbatim}
\newenvironment{reference}{\comment}{\endcomment}
%\newenvironment{reference}{}{}
\newenvironment{slogan}{\comment}{\endcomment}
\newenvironment{history}{\comment}{\endcomment}

% For commutative diagrams we use Xy-pic
\usepackage[all]{xy}

% We use 2cell for 2-commutative diagrams.
\xyoption{2cell}
\UseAllTwocells

% We use multicol for the list of chapters between chapters
\usepackage{multicol}

% This is generally recommended for better output
\usepackage{lmodern}
\usepackage[T1]{fontenc}

% For cross-file-references
\usepackage{xr-hyper}

% Package for hypertext links:
\usepackage{hyperref}

% For any local file, say "hello.tex" you want to link to please
% use \externaldocument[hello-]{hello}
\externaldocument[introduction-]{introduction}
\externaldocument[conventions-]{conventions}
\externaldocument[sets-]{sets}
\externaldocument[categories-]{categories}
\externaldocument[topology-]{topology}
\externaldocument[sheaves-]{sheaves}
\externaldocument[sites-]{sites}
\externaldocument[stacks-]{stacks}
\externaldocument[fields-]{fields}
\externaldocument[algebra-]{algebra}
\externaldocument[brauer-]{brauer}
\externaldocument[homology-]{homology}
\externaldocument[derived-]{derived}
\externaldocument[simplicial-]{simplicial}
\externaldocument[more-algebra-]{more-algebra}
\externaldocument[smoothing-]{smoothing}
\externaldocument[modules-]{modules}
\externaldocument[sites-modules-]{sites-modules}
\externaldocument[injectives-]{injectives}
\externaldocument[cohomology-]{cohomology}
\externaldocument[sites-cohomology-]{sites-cohomology}
\externaldocument[dga-]{dga}
\externaldocument[dpa-]{dpa}
\externaldocument[sdga-]{sdga}
\externaldocument[hypercovering-]{hypercovering}
\externaldocument[schemes-]{schemes}
\externaldocument[constructions-]{constructions}
\externaldocument[properties-]{properties}
\externaldocument[morphisms-]{morphisms}
\externaldocument[coherent-]{coherent}
\externaldocument[divisors-]{divisors}
\externaldocument[limits-]{limits}
\externaldocument[varieties-]{varieties}
\externaldocument[topologies-]{topologies}
\externaldocument[descent-]{descent}
\externaldocument[perfect-]{perfect}
\externaldocument[more-morphisms-]{more-morphisms}
\externaldocument[flat-]{flat}
\externaldocument[groupoids-]{groupoids}
\externaldocument[more-groupoids-]{more-groupoids}
\externaldocument[etale-]{etale}
\externaldocument[chow-]{chow}
\externaldocument[intersection-]{intersection}
\externaldocument[pic-]{pic}
\externaldocument[weil-]{weil}
\externaldocument[adequate-]{adequate}
\externaldocument[dualizing-]{dualizing}
\externaldocument[duality-]{duality}
\externaldocument[discriminant-]{discriminant}
\externaldocument[derham-]{derham}
\externaldocument[local-cohomology-]{local-cohomology}
\externaldocument[algebraization-]{algebraization}
\externaldocument[curves-]{curves}
\externaldocument[resolve-]{resolve}
\externaldocument[models-]{models}
\externaldocument[functors-]{functors}
\externaldocument[equiv-]{equiv}
\externaldocument[pione-]{pione}
\externaldocument[etale-cohomology-]{etale-cohomology}
\externaldocument[proetale-]{proetale}
\externaldocument[relative-cycles-]{relative-cycles}
\externaldocument[more-etale-]{more-etale}
\externaldocument[trace-]{trace}
\externaldocument[crystalline-]{crystalline}
\externaldocument[spaces-]{spaces}
\externaldocument[spaces-properties-]{spaces-properties}
\externaldocument[spaces-morphisms-]{spaces-morphisms}
\externaldocument[decent-spaces-]{decent-spaces}
\externaldocument[spaces-cohomology-]{spaces-cohomology}
\externaldocument[spaces-limits-]{spaces-limits}
\externaldocument[spaces-divisors-]{spaces-divisors}
\externaldocument[spaces-over-fields-]{spaces-over-fields}
\externaldocument[spaces-topologies-]{spaces-topologies}
\externaldocument[spaces-descent-]{spaces-descent}
\externaldocument[spaces-perfect-]{spaces-perfect}
\externaldocument[spaces-more-morphisms-]{spaces-more-morphisms}
\externaldocument[spaces-flat-]{spaces-flat}
\externaldocument[spaces-groupoids-]{spaces-groupoids}
\externaldocument[spaces-more-groupoids-]{spaces-more-groupoids}
\externaldocument[bootstrap-]{bootstrap}
\externaldocument[spaces-pushouts-]{spaces-pushouts}
\externaldocument[spaces-chow-]{spaces-chow}
\externaldocument[groupoids-quotients-]{groupoids-quotients}
\externaldocument[spaces-more-cohomology-]{spaces-more-cohomology}
\externaldocument[spaces-simplicial-]{spaces-simplicial}
\externaldocument[spaces-duality-]{spaces-duality}
\externaldocument[formal-spaces-]{formal-spaces}
\externaldocument[restricted-]{restricted}
\externaldocument[spaces-resolve-]{spaces-resolve}
\externaldocument[formal-defos-]{formal-defos}
\externaldocument[defos-]{defos}
\externaldocument[cotangent-]{cotangent}
\externaldocument[examples-defos-]{examples-defos}
\externaldocument[algebraic-]{algebraic}
\externaldocument[examples-stacks-]{examples-stacks}
\externaldocument[stacks-sheaves-]{stacks-sheaves}
\externaldocument[criteria-]{criteria}
\externaldocument[artin-]{artin}
\externaldocument[quot-]{quot}
\externaldocument[stacks-properties-]{stacks-properties}
\externaldocument[stacks-morphisms-]{stacks-morphisms}
\externaldocument[stacks-limits-]{stacks-limits}
\externaldocument[stacks-cohomology-]{stacks-cohomology}
\externaldocument[stacks-perfect-]{stacks-perfect}
\externaldocument[stacks-introduction-]{stacks-introduction}
\externaldocument[stacks-more-morphisms-]{stacks-more-morphisms}
\externaldocument[stacks-geometry-]{stacks-geometry}
\externaldocument[moduli-]{moduli}
\externaldocument[moduli-curves-]{moduli-curves}
\externaldocument[examples-]{examples}
\externaldocument[exercises-]{exercises}
\externaldocument[guide-]{guide}
\externaldocument[desirables-]{desirables}
\externaldocument[coding-]{coding}
\externaldocument[obsolete-]{obsolete}
\externaldocument[fdl-]{fdl}
\externaldocument[index-]{index}

% Theorem environments.
%
\theoremstyle{plain}
\newtheorem{theorem}[subsection]{Theorem}
\newtheorem{proposition}[subsection]{Proposition}
\newtheorem{lemma}[subsection]{Lemma}

\theoremstyle{definition}
\newtheorem{definition}[subsection]{Definition}
\newtheorem{example}[subsection]{Example}
\newtheorem{exercise}[subsection]{Exercise}
\newtheorem{situation}[subsection]{Situation}

\theoremstyle{remark}
\newtheorem{remark}[subsection]{Remark}
\newtheorem{remarks}[subsection]{Remarks}

\numberwithin{equation}{subsection}

% Macros
%
\def\lim{\mathop{\mathrm{lim}}\nolimits}
\def\colim{\mathop{\mathrm{colim}}\nolimits}
\def\Spec{\mathop{\mathrm{Spec}}}
\def\Hom{\mathop{\mathrm{Hom}}\nolimits}
\def\Ext{\mathop{\mathrm{Ext}}\nolimits}
\def\SheafHom{\mathop{\mathcal{H}\!\mathit{om}}\nolimits}
\def\SheafExt{\mathop{\mathcal{E}\!\mathit{xt}}\nolimits}
\def\Sch{\mathit{Sch}}
\def\Mor{\mathop{\mathrm{Mor}}\nolimits}
\def\Ob{\mathop{\mathrm{Ob}}\nolimits}
\def\Sh{\mathop{\mathit{Sh}}\nolimits}
\def\NL{\mathop{N\!L}\nolimits}
\def\CH{\mathop{\mathrm{CH}}\nolimits}
\def\proetale{{pro\text{-}\acute{e}tale}}
\def\etale{{\acute{e}tale}}
\def\QCoh{\mathit{QCoh}}
\def\Ker{\mathop{\mathrm{Ker}}}
\def\Im{\mathop{\mathrm{Im}}}
\def\Coker{\mathop{\mathrm{Coker}}}
\def\Coim{\mathop{\mathrm{Coim}}}

% Boxtimes
%
\DeclareMathSymbol{\boxtimes}{\mathbin}{AMSa}{"02}

%
% Macros for moduli stacks/spaces
%
\def\QCohstack{\mathcal{QC}\!\mathit{oh}}
\def\Cohstack{\mathcal{C}\!\mathit{oh}}
\def\Spacesstack{\mathcal{S}\!\mathit{paces}}
\def\Quotfunctor{\mathrm{Quot}}
\def\Hilbfunctor{\mathrm{Hilb}}
\def\Curvesstack{\mathcal{C}\!\mathit{urves}}
\def\Polarizedstack{\mathcal{P}\!\mathit{olarized}}
\def\Complexesstack{\mathcal{C}\!\mathit{omplexes}}
% \Pic is the operator that assigns to X its picard group, usage \Pic(X)
% \Picardstack_{X/B} denotes the Picard stack of X over B
% \Picardfunctor_{X/B} denotes the Picard functor of X over B
\def\Pic{\mathop{\mathrm{Pic}}\nolimits}
\def\Picardstack{\mathcal{P}\!\mathit{ic}}
\def\Picardfunctor{\mathrm{Pic}}
\def\Deformationcategory{\mathcal{D}\!\mathit{ef}}
\usepackage{cleveref}
\crefname{originallemma}{original lemma}{original lemmas}
\Crefname{originallemma}{Original lemma}{Original lemmas}
\makeatletter
\newlabel{modules-lemma-i-star-exact}{{01AX}{source}{}{sourcepointer-1}{}}
\newlabel{modules-lemma-i-star-exact@cref}{{[originallemma][0][]01AX}{[1][0][]source}{}{}{}}
\makeatother
\crefname{originallemma}{original lemma}{original lemmas}
\Crefname{originallemma}{Original lemma}{Original lemmas}
\makeatletter
\newlabel{cohomology-lemma-finite-cohomology}{{08DN}{source}{}{sourcepointer-2}{}}
\newlabel{cohomology-lemma-finite-cohomology@cref}{{[originallemma][0][]08DN}{[1][0][]source}{}{}{}}
\makeatother
\crefname{originallemma}{original lemma}{original lemmas}
\Crefname{originallemma}{Original lemma}{Original lemmas}
\makeatletter
\newlabel{cohomology-lemma-cone-pseudo-coherent}{{08CD}{source}{}{sourcepointer-3}{}}
\newlabel{cohomology-lemma-cone-pseudo-coherent@cref}{{[originallemma][0][]08CD}{[1][0][]source}{}{}{}}
\makeatother
\title{Pseudo-coherence is preserved and reflected by the pushforward along a closed immersion}
\author{The Stacks Project Authors}
\date{stacks-project:perfect.tex:local-snapshot-20261012; GNU FDL 1.2 or later}
\begin{document}
\maketitle
\setcounter{subsection}{1}
\noindent\textit{Editorial note.} All delivered mathematical text is the author's own, copied from the pinned Stacks Project chapter file perfect.tex; no statement or proof step was rewritten, reordered or removed. The counted span is the author environment and its own proof environment. References inside the span are handed to the assembler as navigation stubs that carry the canonical tag of the cited tag only; whether each of them is load-bearing is recorded on the stub itself.\par\medskip
\noindent\textit{Allowed standard prerequisites (editorial).}\begin{itemize}
\item The pushforward along a closed immersion of ringed spaces is exact, so the derived and the plain pushforward agree. used at proof lines 3151-3153. Lets the proof work with the plain pushforward throughout.
\item A finitely generated cohomology module of a pseudo-coherent complex admits, locally, a surjection from a finite free module onto it. invoked at proof lines 3160-3164. Produces the map from a shifted free module onto the top cohomology group.
\item In a distinguished triangle, if two terms are pseudo-coherent then so is the third, so cutting off the top cohomology preserves the property. invoked at proof lines 3166-3168 and again at proof lines 3170-3175. Gives that the cone is pseudo-coherent and then that the pushed forward complex is pseudo-coherent.
\end{itemize}\medskip
\ifdefined\hypertarget\hypertarget{sourcepointer-1}{}\fi
\noindent\textit{Source reference.} Original lemma 01AX: sources/stacks-chow/tags.txt, source lines 1--1. Cited at proof lines 3151-3153 as the fact used throughout that the derived pushforward is the plain one; all subsequent arguments are inside the counted span.\par\medskip
\ifdefined\hypertarget\hypertarget{sourcepointer-2}{}\fi
\noindent\textit{Source reference.} Original lemma 08DN: sources/stacks-chow/tags.txt, source lines 1--1. Invoked at proof lines 3160-3162 to see that the top cohomology module is of finite type, which produces the surjection from a free module; the construction of the triangle that follows is inside the counted span.\par\medskip
\ifdefined\hypertarget\hypertarget{sourcepointer-3}{}\fi
\noindent\textit{Source reference.} Original lemma 08CD: sources/stacks-chow/tags.txt, source lines 1--1. Applied twice, at proof lines 3166-3168 to see that the cone is pseudo-coherent and at proof lines 3172-3176 to transfer the property along the pushed forward triangle; both applications are inside the counted span.\par\medskip
\begin{lemma}
\label{lemma-closed-push-pseudo-coherent}
Let $i : Z \to X$ be a morphism of ringed spaces such that
$i$ is a closed immersion of underlying topological spaces and such that
$i_*\mathcal{O}_Z$ is pseudo-coherent as an $\mathcal{O}_X$-module.
Let $E \in D(\mathcal{O}_Z)$. Then $E$ is $m$-pseudo-coherent
if and only if $Ri_*E$ is $m$-pseudo-coherent.
\end{lemma}

\begin{proof}
Throughout this proof we will use that $i_*$ is an exact functor, and
hence that $Ri_* = i_*$, see Modules, Lemma \ref{modules-lemma-i-star-exact}.

\medskip\noindent
Assume $E$ is $m$-pseudo-coherent. Let $x \in X$. We will find a neighbourhood
of $x$ such that $i_*E$ is $m$-pseudo-coherent on it. If $x \not \in Z$
then this is clear. Thus we may assume $x \in Z$. We will use
that $U \cap Z$ for $x \in U \subset X$ open form a fundamental system of
neighbourhoods of $x$ in $Z$. After shrinking $X$ we may assume $E$ is
bounded above. We will argue by induction on
the largest integer $p$ such that $H^p(E)$ is nonzero. If $p < m$, then
there is nothing to prove. If $p \geq m$, then $H^p(E)$ is an
$\mathcal{O}_Z$-module of finite type, see
Cohomology, Lemma \ref{cohomology-lemma-finite-cohomology}.
Thus we may choose, after shrinking $X$, a map
$\mathcal{O}_Z^{\oplus n}[-p] \to E$ which induces a surjection
$\mathcal{O}_Z^{\oplus n} \to H^p(E)$. Choose a distinguished triangle
$$
\mathcal{O}_Z^{\oplus n}[-p] \to E \to C \to \mathcal{O}_Z^{\oplus n}[-p + 1]
$$
We see that $H^j(C) = 0$ for $j \geq p$ and that $C$ is $m$-pseudo-coherent
by Cohomology, Lemma \ref{cohomology-lemma-cone-pseudo-coherent}.
By induction we see that $i_*C$ is $m$-pseudo-coherent on $X$.
Since $i_*\mathcal{O}_Z$ is $m$-pseudo-coherent on $X$ as well, we conclude
from the distinguished triangle
$$
i_*\mathcal{O}_Z^{\oplus n}[-p] \to i_*E \to i_*C \to
i_*\mathcal{O}_Z^{\oplus n}[-p + 1]
$$
and 
Cohomology, Lemma \ref{cohomology-lemma-cone-pseudo-coherent}
that $i_*E$ is $m$-pseudo-coherent.

\medskip\noindent
Assume that $i_*E$ is $m$-pseudo-coherent. Let $z \in Z$.
We will find a neighbourhood of $z$ such that $E$
is $m$-pseudo-coherent on it. We will use
that $U \cap Z$ for $z \in U \subset X$ open form a fundamental system of
neighbourhoods of $z$ in $Z$. After shrinking $X$ we may assume $i_*E$
and hence $E$ is bounded above. We will argue by induction on
the largest integer $p$ such that $H^p(E)$ is nonzero. If $p < m$, then
there is nothing to prove. If $p \geq m$, then $H^p(i_*E) = i_*H^p(E)$
is an $\mathcal{O}_X$-module of finite type, see
Cohomology, Lemma \ref{cohomology-lemma-finite-cohomology}.
Choose a complex $\mathcal{E}^\bullet$ of $\mathcal{O}_Z$-modules
representing $E$. We may choose, after shrinking $X$,
a map $\alpha : \mathcal{O}_X^{\oplus n}[-p] \to i_*\mathcal{E}^\bullet$
which induces a surjection
$\mathcal{O}_X^{\oplus n} \to i_*H^p(\mathcal{E}^\bullet)$.
By adjunction we find a map
$\alpha : \mathcal{O}_Z^{\oplus n}[-p] \to \mathcal{E}^\bullet$
which induces a surjection
$\mathcal{O}_Z^{\oplus n} \to H^p(\mathcal{E}^\bullet)$.
Choose a distinguished triangle
$$
\mathcal{O}_Z^{\oplus n}[-p] \to E \to C \to \mathcal{O}_Z^{\oplus n}[-p + 1]
$$
We see that $H^j(C) = 0$ for $j \geq p$. From the distinguished triangle
$$
i_*\mathcal{O}_Z^{\oplus n}[-p] \to i_*E \to i_*C \to
i_*\mathcal{O}_Z^{\oplus n}[-p + 1]
$$
the fact that $i_*\mathcal{O}_Z$ is pseudo-coherent
and 
Cohomology, Lemma \ref{cohomology-lemma-cone-pseudo-coherent}
we conclude that $i_*C$ is $m$-pseudo-coherent.
By induction we conclude that $C$ is $m$-pseudo-coherent.
By Cohomology, Lemma \ref{cohomology-lemma-cone-pseudo-coherent}
again we conclude that $E$ is $m$-pseudo-coherent.
\end{proof}

\end{document}
